% Generated by tools/build_new_dists_anthology_sections.py; do not edit by hand.

\section{Braid, TQFT, and quantum-image groups}\label{proof:new-dist:braid_tqft}

\resultbox{\textbf{Canonical testing artifacts.}\par\smallskip Exact backend: \path{new_dists/braid_tqft_quantum_images_sigma_mgfs/verification.py}.\par Regression: \path{new_dists/braid_tqft_quantum_images_sigma_mgfs/test_verification.py}.\par Marimo: \path{marimo_notebooks/braid_tqft.py}.}

\subsection*{Scope and outcome}

A finite braid or TQFT image has a uniform measure; an infinite image does
not, and its canonical unitary substitute is normalized Haar measure on the
compact closure.  We prove the generalized Molien/class-power theorem in both
settings and test each mathematically distinct mechanism on a specified
representation.  The program constructs a $192$-matrix Ising braid lift, the
$27$-matrix qutrit Heisenberg group, a $96$-matrix rank-one Weil lift, the
$24$-element projective qubit Clifford group, and regular and affine
phase-space permutation groups through degree $25$.  All 35 finite formula
comparisons and all 1,440 affine fixed-power checks pass.  Sixteen seeded Haar
diagnostics corroborate the exact $U$, $SU$, $O$, and $USp$ conclusions.  The
generic Jones, mapping-class, and Hadamard statements are recorded with their
necessary image/closure hypotheses; no unspecified representation is treated
as though it determined a matrix group.

\subsection{The probability distinction and the common target}

Let $V$ be a $d$-dimensional complex unitary module.  For a finite represented
image $I\leq U(V)$, define
\begin{equation}\label{nd:braid_tqft:eq:finite-def}
\MGF_{I,V}(\mathbf t)
=\frac1{|I|}\sum_{x\in I}\prod_{a\geq1}
\det(I-t_a x)^{-1}.
\end{equation}
If a unitary representation has infinite image, there is no uniform
probability measure on that countable image.  Put instead
\[
K=\overline{\rho(\Gamma)}\leq U(V)
\]
and use normalized Haar measure:
\begin{equation}\label{nd:braid_tqft:eq:compact-def}
\MGF^{\mathrm{cl}}_{\rho,V}(\mathbf t)
=\int_K\prod_{a\geq1}\det(I-t_a k)^{-1}\,dk.
\end{equation}
This agrees with \eqref{nd:braid_tqft:eq:finite-def} for finite $K$ and is unchanged when a
dense subgroup is replaced by its closure.  Thus ``braid representation'' or
``TQFT representation'' is not enough data: one must specify the represented
image, a finite central lift when the action is projective, or the compact
closure.

For a partition or composition $\tau=(\tau_1,\ldots,\tau_s)$, set
\[
W_\tau(V)=\bigotimes_{j=1}^s\Sym^{\tau_j}V,
\qquad |\tau|=\sum_j\tau_j.
\]
The coefficient tested throughout is
\[
[m_\tau]\MGF_{G,V}=\dim W_\tau(V)^G.
\]

\subsection{Universal finite-image and compact-closure theorem}

\begin{theorem}[Generalized Molien/class-power identity]\label{nd:braid_tqft:thm:universal}
For a finite unitary image $I$ with character $\chi_V$,
\begin{align}
\MGF_{I,V}
&=\frac1{|I|}\sum_{x\in I}
\exp\!\left(\sum_{r\geq1}\frac{p_r}{r}\chi_V(x^r)\right)\label{nd:braid_tqft:eq:element-power}\\
&=\sum_{[x]\subseteq I}\frac1{|C_I(x)|}
\exp\!\left(\sum_{r\geq1}\frac{p_r}{r}\chi_V(x^r)\right),\label{nd:braid_tqft:eq:class-power}
\end{align}
where $p_r=\sum_a t_a^r$.  Moreover
\begin{equation}\label{nd:braid_tqft:eq:coeff}
\boxed{[m_\tau]\MGF_{I,V}=\dim W_\tau(V)^I.}
\end{equation}
The same statements hold for compact $K$ after replacing the finite average
by normalized Haar integration.
\end{theorem}

\begin{proof}
If the eigenvalues of a matrix $A$ are $\lambda_1,\ldots,\lambda_d$, then as a
formal power-series identity
\[
\det(I-tA)^{-1}
=\prod_j(1-t\lambda_j)^{-1}
=\exp\!\left(\sum_{r\geq1}\frac{t^r}{r}\Tr(A^r)\right).
\]
Multiply over $t_a$, substitute $A=\rho(x)$, and use
$\Tr(\rho(x)^r)=\chi_V(x^r)$ to obtain \eqref{nd:braid_tqft:eq:element-power}.
Grouping a finite sum into conjugacy classes uses
$|[x]|/|I|=1/|C_I(x)|$ and gives \eqref{nd:braid_tqft:eq:class-power}.

The coefficient of $t^n$ in $\det(I-tA)^{-1}$ is
$\Tr(\Sym^n A)$.  Hence the coefficient indexed by $\tau$ is the average of
the character of $W_\tau(V)$.  The averaging operator
\[
P_\tau=\frac1{|I|}\sum_{x\in I}\rho_{W_\tau}(x)
\]
is the Reynolds projection onto $W_\tau(V)^I$, so its trace equals the
dimension in \eqref{nd:braid_tqft:eq:coeff}.  For compact $K$, Haar invariance gives the same
idempotent after replacing the sum by an integral.
\end{proof}

The theorem shows that the exact finite answer is determined by conjugacy
classes, power maps, and the represented character.  It does \emph{not} say
that the abstract source group or a TQFT label determines those data.

\subsubsection{Projective representations and scalar selection}

Suppose a projective image becomes honest on a finite central extension
\[
1\longrightarrow Z\longrightarrow\widetilde I\longrightarrow I
\longrightarrow1.
\]
Theorem~\ref{nd:braid_tqft:thm:universal} applies to $\widetilde I$.  If $z\in Z$ acts as
$\xi(z)I$, then $z$ acts on $W_\tau(V)$ as $\xi(z)^{|\tau|}$.  Therefore
\begin{equation}\label{nd:braid_tqft:eq:central-rule}
[m_\tau]\MGF_{\widetilde I,V}=0
\quad\text{unless}\quad
\xi(z)^{|\tau|}=1\ \text{for every }z\in Z.
\end{equation}
This is a necessary condition, not a promise of a nonzero invariant in every
allowed degree.  It also explains why the one-sided series depends on the
chosen finite lift.

\subsection{Finite Ising braid image}

Let
\[
H=\frac1{\sqrt2}\begin{pmatrix}1&1\\1&-1\end{pmatrix},\qquad
R=e^{-\pi i/8}\begin{pmatrix}1&0\\0&i\end{pmatrix},
\]
and put
\begin{equation}\label{nd:braid_tqft:eq:ising-generators}
\rho(\sigma_1)=R,\qquad \rho(\sigma_2)=HRH.
\end{equation}
Direct multiplication gives
\[
\rho(\sigma_1)\rho(\sigma_2)\rho(\sigma_1)
=\rho(\sigma_2)\rho(\sigma_1)\rho(\sigma_2),
\]
so \eqref{nd:braid_tqft:eq:ising-generators} is an honest two-dimensional representation of
$B_3$.  The verifier starts from these matrices and closes under
multiplication.  It obtains $192$ matrices, $24$ matrices modulo scalar phase,
and exactly eight scalar matrices $\mu_8I$.  Thus \eqref{nd:braid_tqft:eq:central-rule}
forces $8\mid |\tau|$.

For each matrix $A$, the direct route diagonalizes $A$ and expands
$\det(I-tA)^{-1}$.  The independent power route computes $\Tr(A^r)$ by matrix
multiplication and uses Newton's recurrence
\begin{equation}\label{nd:braid_tqft:eq:newton}
h_0(A)=1,\qquad h_n(A)=\frac1n\sum_{r=1}^n
\Tr(A^r)h_{n-r}(A).
\end{equation}
The literal matrix averages agree:
\begin{center}
\begin{tabular}{@{}rrrrrrrr@{}}
\toprule
$\tau$ & $(1)$ & $(2)$ & $(4)$ & $(8)$ & $(7,1)$ & $(4,4)$ & $(1^8)$\\
\midrule
direct matrices &0&0&0&0&1&1&7\\
power formula   &0&0&0&0&1&1&7\\
\bottomrule
\end{tabular}
\end{center}
At $(t_1,t_2,t_3)=(0.023,-0.031,0.041)$, direct determinants and the
truncated power-character exponential agree to $2.2\times10^{-18}$; the
braid-relation residual is $3.0\times10^{-16}$.  The power series is truncated
at $r=60$ only for this nonzero numerical evaluation.  All displayed
coefficients use the exact finite recurrence \eqref{nd:braid_tqft:eq:newton} through the
required degree.

\begin{remark}
This example validates a particular finite Ising/Clifford braid image.  It
does not imply that generic Jones, Hecke, Temperley--Lieb, or BMW sectors have
finite image.  Density results for broad Jones--Wenzl families show why the
finite/closure distinction is substantive \cite{nd:braid_tqft:flw}.
\end{remark}

\subsection{Monomial finite image: the qutrit Heisenberg group}

Suppose a represented finite group is monomial in a chosen basis.  Write
$x=D(\alpha)P_\pi$.  For a cycle $C=(j_1\cdots j_\ell)$ of $\pi$, define
$\alpha_C=\prod_{j\in C}\alpha_j$.  Reordering the basis cycle by cycle gives
\begin{equation}\label{nd:braid_tqft:eq:monomial-det}
\det(I-tx)=\prod_{C\in\operatorname{Cyc}(\pi)}
(1-\alpha_Ct^{|C|}),
\end{equation}
and hence
\begin{equation}\label{nd:braid_tqft:eq:monomial-mgf}
\boxed{\MGF_{I,V}=\frac1{|I|}\sum_{x\in I}
\prod_{a,C}(1-\alpha_C(x)t_a^{|C|})^{-1}.}
\end{equation}

The concrete test is the qutrit Heisenberg group
\[
H_3=\{\omega^cX^aZ^b:a,b,c\in\F_3\},\qquad
Xe_j=e_{j+1},\quad Ze_j=\omega^je_j,
\]
where $\omega=e^{2\pi i/3}$.  All $27$ dense $3\times3$ matrices are built.
The code separately reads the nonzero entry in each column, traverses the
permutation cycles, multiplies their phases, and evaluates
\eqref{nd:braid_tqft:eq:monomial-mgf}.  The dense and cycle-phase multivariate averages
agree to $1.8\times10^{-20}$.

Here the central phase rule and the two spectral sectors simplify to
\begin{equation}\label{nd:braid_tqft:eq:h3formula}
[m_\tau]\MGF_{H_3,V}=
\begin{cases}
\displaystyle \frac19\prod_j\binom{\tau_j+2}{\tau_j}
+\frac89\prod_j\mathbf1_{3\mid\tau_j},&3\mid|\tau|,\\[6pt]
0,&3\nmid|\tau|.
\end{cases}
\end{equation}
Indeed, the three central matrices contribute the first term after the phase
filter, while every noncentral matrix has spectrum a scalar multiple of all
three cube roots and contributes the second term.

\begin{center}
\begin{tabular}{@{}lrrrrrrr@{}}
\toprule
$\tau$&$(1)$&$(2)$&$(3)$&$(2,1)$&$(4)$&$(3,3)$&$(6)$\\
\midrule
direct matrices&0&0&2&2&0&12&4\\
formula \eqref{nd:braid_tqft:eq:h3formula}&0&0&2&2&0&12&4\\
\bottomrule
\end{tabular}
\end{center}
This is a representative finite Gaussian/Heisenberg core of the type that
occurs in metaplectic and parafermionic constructions.  Identifying a
particular physical braid sector with this group requires its own generators;
finiteness alone does not imply a universal asymptotic law.

\subsection{Rank-one Weil lift and the phase issue}

Over $\F_3$, define the Fourier and chirp matrices
\begin{equation}\label{nd:braid_tqft:eq:weil-generators}
F_{xy}=3^{-1/2}\omega^{xy},\qquad
P_{xx}=\omega^{x^2/2},
\end{equation}
where $1/2=2$ in $\F_3$.  Simultaneously track the generators
\[
S=\begin{pmatrix}0&-1\\1&0\end{pmatrix},\qquad
T=\begin{pmatrix}1&0\\1&1\end{pmatrix}
\quad\text{in }SL_2(\F_3).
\]
Projective matrix closure gives all $24$ elements of $SL_2(\F_3)$; retaining
the generated scalar phases gives a $96$-matrix lift.

The finite-field Weil character has the form
\begin{equation}\label{nd:braid_tqft:eq:weil-character}
\chi_\omega(\widetilde A)=\gamma(\widetilde A)
|\ker(A-I)|^{1/2},\qquad |\gamma(\widetilde A)|=1,
\end{equation}
with a lift-dependent Weil-index phase \cite{nd:braid_tqft:thomas}.  Consequently
\begin{equation}\label{nd:braid_tqft:eq:weil-magnitude}
|\chi_\omega(\widetilde A)|^2=|\ker(A-I)|.
\end{equation}
The verifier checks \eqref{nd:braid_tqft:eq:weil-magnitude} for every $A\in SL_2(\F_3)$.
The distribution is
\begin{center}
\begin{tabular}{@{}crrr@{}}
\toprule
$|\ker(A-I)|$&1&3&9\\
number of matrices&15&8&1\\
observed $|\Tr\omega(A)|^2$&1&3&9\\
\bottomrule
\end{tabular}
\end{center}

For the full lift, the phase-sensitive class-power formula is
\begin{equation}\label{nd:braid_tqft:eq:weil-class}
\MGF^{\mathrm{Weil}}=
\sum_{[\widetilde A]}\frac1{|C(\widetilde A)|}
\exp\!\left(\sum_{r\geq1}\frac{p_r}{r}
\chi_\omega(\widetilde A^r)\right).
\end{equation}
The direct and power-character coefficient averages give
\begin{center}
\begin{tabular}{@{}lrrrrr@{}}
\toprule
$\tau$&$(1)$&$(2)$&$(4)$&$(3,1)$&$(4,4)$\\
\midrule
direct lift&0&0&1&2&6\\
formula \eqref{nd:braid_tqft:eq:weil-class}&0&0&1&2&6\\
\bottomrule
\end{tabular}
\end{center}
The full numerical MGF error is $2.3\times10^{-16}$.

Equation \eqref{nd:braid_tqft:eq:weil-magnitude} controls balanced trace moments, but it does
not determine the holomorphic series \eqref{nd:braid_tqft:eq:weil-class}: the latter also
needs the joint phases of $A,A^2,\ldots$.  This is exactly why fixed-space
magnitude stabilization cannot by itself prove convergence of the full Weil
sigma--MGF.

The same rank-one construction may be viewed as a finite quotient of a torus
mapping-class representation.  Higher-genus or composite-level formulas
remain exact class sums only after the relevant lift and Weil character have
been specified; Chinese-remainder structure alone does not remove the local
phase data.

\subsection{Balanced series and unitary designs}

Scalar phases often annihilate the one-sided series.  The appropriate
design-sensitive object is
\begin{equation}\label{nd:braid_tqft:eq:balanced}
\BMGF_{G,V}(\mathbf s,\mathbf t)=\frac1{|G|}\sum_{g\in G}
\prod_i\det(I-s_ig)^{-1}\prod_j\det(I-t_j\overline g)^{-1}.
\end{equation}
Character expansion and Reynolds averaging prove
\begin{equation}\label{nd:braid_tqft:eq:balanced-coeff}
[m_\alpha(\mathbf s)m_\beta(\mathbf t)]\BMGF_{G,V}
=\dim\left[
\left(\bigotimes_i\Sym^{\alpha_i}V\right)\otimes
\left(\bigotimes_j\Sym^{\beta_j}V^*\right)
\right]^G.
\end{equation}
Only $|\alpha|=|\beta|$ is intrinsically projective.  For
$\alpha=\beta=(1^k)$, the coefficient becomes the frame potential
\begin{equation}\label{nd:braid_tqft:eq:frame}
\Phi_k(G)=\frac1{|G|}\sum_{g\in G}|\Tr g|^{2k}
=\dim\operatorname{End}_G(V^{\otimes k}).
\end{equation}

If $G$ is a unitary $t$-design, every balanced coefficient with
$|\alpha|,|\beta|\leq t$ equals its Haar $U(d)$ value.  Conversely, equality
of \eqref{nd:braid_tqft:eq:frame} with Haar at degree $k$ forces equality of the nested
invariant spaces and hence equality of the $k$th moment projector.  By
Schur--Weyl duality,
\[
\Phi_k(U(d))=\sum_{\substack{\lambda\vdash k\\\ell(\lambda)\leq d}}
(f^\lambda)^2.
\]

Closing $H$ and $S=\operatorname{diag}(1,i)$ modulo scalar phase gives the
$24$-element projective one-qubit Clifford group.  Direct matrix averaging
and the hook-length formula give
\begin{center}
\begin{tabular}{@{}crrrr@{}}
\toprule
$k$&1&2&3&4\\
\midrule
$\Phi_k(\mathrm{Cliff}_1)$&1&2&5&15\\
$\Phi_k(U(2))$&1&2&5&14\\
\bottomrule
\end{tabular}
\end{center}
Thus the tested group agrees through degree three and first differs at degree
four, consistent with the general multiqubit design theorem and its extra
fourth-order stabilizer invariant \cite{nd:braid_tqft:clifford4}.

\subsection{Weyl--Heisenberg regular and affine phase-space actions}

\subsubsection{Regular action and the prime SIC formula}

Let $A_d\cong(\mathbb Z/d\mathbb Z)^2$ act regularly on itself, viewed as the
$d^2$ projective outcomes in a Weyl--Heisenberg covariant SIC setting.  An
element $a$ of order $o(a)$ has $d^2/o(a)$ cycles, so
\begin{equation}\label{nd:braid_tqft:eq:regular-general}
\MGF_{A_d,\mathrm{Reg}}=
\frac1{d^2}\sum_{a\in A_d}\prod_i
(1-t_i^{o(a)})^{-d^2/o(a)}.
\end{equation}
For prime $p$, every nonidentity element has order $p$:
\begin{equation}\label{nd:braid_tqft:eq:regular-prime}
\boxed{\MGF_{A_p,\mathrm{Reg}}=
\frac1{p^2}\left[
\prod_i(1-t_i)^{-p^2}+(p^2-1)\prod_i(1-t_i^p)^{-p}
\right].}
\end{equation}
In particular
\begin{equation}\label{nd:braid_tqft:eq:regular-pair}
[m_{(1,1)}]\MGF_{A_p,\mathrm{Reg}}=p^2,
\end{equation}
because only the identity fixes a point and the pair coefficient is the
average squared fixed-point count.

The verifier constructs all permutation matrices for $p=2,3,5$.  Selected
exact comparisons are
\begin{center}
\begin{tabular}{@{}crrrr@{}}
\toprule
$p$&permutation degree&$[m_{11}]$&$[m_2]$&$[m_{p,p}]$\\
\midrule
2&4&4&4&28\\
3&9&9&5&3,033\\
5&25&25&13&564,110,025\\
\bottomrule
\end{tabular}
\end{center}
Every entry agrees with coefficient extraction from \eqref{nd:braid_tqft:eq:regular-prime};
full nonzero determinant evaluations agree to $2.3\times10^{-16}$.  The growth
in \eqref{nd:braid_tqft:eq:regular-pair} proves that the raw regular-action series has no
coefficientwise $p\to\infty$ limit.  The contextual link between prime
Weyl--Heisenberg SICs and Clifford normalizers is discussed in \cite{nd:braid_tqft:zhu-sic}.

\subsubsection{Affine symmetry groups}

Let $\Gamma=A\rtimes K$ act on the additive space $A$.  Write
$g=(b,A_0)$ with action $x\mapsto A_0x+b$ and put
\[
b_r=(I+A_0+\cdots+A_0^{r-1})b.
\]
Then
\[
g^r(x)=A_0^rx+b_r,
\]
so the fixed-point equation is $(I-A_0^r)x=b_r$.  Therefore
\begin{equation}\label{nd:braid_tqft:eq:affine-fixed}
\boxed{
|\Fix(g^r)|=
\begin{cases}
|\ker(I-A_0^r)|,&b_r\in\operatorname{im}(I-A_0^r),\\
0,&\text{otherwise}.
\end{cases}}
\end{equation}
Fixed points of powers recover cycle counts by M\"obius inversion, hence
\eqref{nd:braid_tqft:eq:affine-fixed} determines the full permutation sigma--MGF.

Ordered-pair orbits are classified by the difference of the two points, so
\begin{equation}\label{nd:braid_tqft:eq:affine-pair}
[m_{(1,1)}]\MGF_{A\rtimes K,A}
=\#(K\backslash A)=1+\#(K\backslash(A\setminus\{0\})).
\end{equation}
For $K=SL_2(\F_p)$, the nonzero vectors form one orbit and the answer is $2$.
The code constructs the full affine permutation groups for $p=2,3$:
\begin{center}
\begin{tabular}{@{}crrrr@{}}
\toprule
$p$&degree&group order&fixed-power checks&pair coefficient\\
\midrule
2&4&24&144&2\\
3&9&216&1,296&2\\
\bottomrule
\end{tabular}
\end{center}
All $1,440$ comparisons of direct fixed points with
\eqref{nd:braid_tqft:eq:affine-fixed}, for powers $1\leq r\leq6$, pass exactly.

\subsection{Compact closures and stable classical laws}

If the represented image is infinite, Theorem~\ref{nd:braid_tqft:thm:universal} applies to
its compact closure $K$.  Four cases must be distinguished.

\subsubsection{Unitary and special-unitary closure}

If $K=U(d)$ in the defining representation, the scalar circle acts on every
positive-degree $W_\tau(V)$ with weight $|\tau|$.  Hence
\begin{equation}\label{nd:braid_tqft:eq:u-trivial}
\boxed{\MGF_{U(d),\C^d}=1.}
\end{equation}
For $SU(d)$, its center $\mu_dI$ implies that a coefficient vanishes unless
$d\mid|\tau|$.  Consequently, for every fixed positive $\tau$, the
coefficient is exactly zero once $d>|\tau|$, and
\begin{equation}\label{nd:braid_tqft:eq:su-limit}
\MGF_{SU(d),\C^d}\longrightarrow1
\quad\text{coefficientwise.}
\end{equation}
The divisibility condition is only necessary.  For example,
$[m_{(1^d)}]=1$ because the determinant tensor is invariant, whereas
$[m_{(d)}]=0$ for $d>1$ in the defining module.

\subsubsection{Orthogonal and compact symplectic closure}

The first fundamental theorems for $O(d)$ and $USp(2n)$ say that stable tensor
invariants are generated by pair contractions with, respectively, a
symmetric or alternating bilinear form.  Repackaging these contractions by
symmetric powers yields the stable identities
\begin{equation}\label{nd:braid_tqft:eq:stable-classical}
\boxed{
\MGF_{O(\infty)}=\ExpS(h_2),\qquad
\MGF_{USp(\infty)}=\ExpS(e_2).}
\end{equation}
This is coefficientwise stability, not equality in every degree at every
finite rank; contraction diagrams can acquire relations outside the stable
range.

The proof of \eqref{nd:braid_tqft:eq:u-trivial}--\eqref{nd:braid_tqft:eq:stable-classical} is algebraic.
As an independent diagnostic, the notebook draws seeded Haar matrices by
complex Gaussian QR for $U$, determinant-normalized QR for $SU$, real QR for
$O$, and quaternionic Gram--Schmidt for $USp$.  It averages the same
symmetric-power characters used in the finite tests.  Representative rows
from $3,000$ samples each are:
\begin{center}
\begin{tabular}{@{}llrrr@{}}
\toprule
group&$\tau$&estimate&target&standard error\\
\midrule
$U(4)$&$(1,1)$&$-0.0279-0.0160i$&0&0.0258\\
$SU(2)$&$(1,1)$&1.0064&1&0.0182\\
$SU(3)$&$(1,1,1)$&$0.9930-0.0013i$&1&0.0412\\
$O(6)$&$(2)$&0.9956&1&0.0182\\
$O(6)$&$(1,1,1,1)$&2.7695&3&0.1503\\
$O(6)$&$(2,2)$&2.1018&2&0.0731\\
$USp(4)$&$(2)$&$-0.0209$&0&0.0178\\
$USp(4)$&$(1,1,1,1)$&3.1935&3&0.1554\\
$USp(4)$&$(2,2)$&1.0291&1&0.0437\\
\bottomrule
\end{tabular}
\end{center}
All sixteen diagnostics lie within the stated five-standard-error/absolute
tolerance.  Worst unitarity residuals are below $6.1\times10^{-15}$.  These
samples test the matrix construction and low-degree consequences; they are
not substituted for the Reynolds and invariant-theory proofs.

\subsection{How the remaining family labels fit}

The universal theorem organizes the broader claims without assigning a
formula to an unspecified representation.

\begin{itemize}
\item \textbf{Jones, Hecke, Temperley--Lieb, and BMW sectors.}  For a specified
finite image, use \eqref{nd:braid_tqft:eq:class-power}.  For an infinite image, use
\eqref{nd:braid_tqft:eq:compact-def}.  If the closure is eventually $U(d_n)$ or $SU(d_n)$,
the one-sided coefficientwise limit is $1$; orthogonal and symplectic closures
give \eqref{nd:braid_tqft:eq:stable-classical}.  Fusion-space dimension growth alone does
not determine the closure.
\item \textbf{Property (F) and finite mapping-class images.}  Property (F)
supplies finiteness, not a universal class table.  One still needs the lift,
power maps, and represented character.  Twisted quantum-double and
metaplectic examples often allow monomial or Weil compression; the
Heisenberg and rank-one Weil sections demonstrate the two reductions.
\item \textbf{Dense mapping-class images.}  Projective density must first be
translated to the honest compact closure relevant to the chosen lift.  Once
that closure and its central scalars are identified, the compact theorem
applies.
\item \textbf{Hessian and Heisenberg-normalizer groups.}  The defining
unitary action uses a Weil--Heisenberg class sum; the phase-space action uses
\eqref{nd:braid_tqft:eq:affine-fixed}.  Their asymptotics depend on fixed-space phases or on
the orbit structure of the complement.
\item \textbf{Hadamard-generated groups.}  Hadamard generation alone does not
imply finiteness.  In the finite qubit example, $H$ and $HS$ are both complex
Hadamard and generate the same lift as $H,S$, since $H(HS)=S$.  Infinite
examples require the closure-Haar formula.
\end{itemize}

The general obstruction to a raw coefficientwise limit is already visible at
balanced fourth tensor degree:
\[
\dim\operatorname{End}_{I_n}(V_n^{\otimes2}).
\]
If this diverges, the corresponding coefficient diverges.  Conversely,
bounded low moments do not determine the full holomorphic series when
power-character phases fail to stabilize.

\subsection{Verification design and reproducibility}

The program uses structurally different routes:
\begin{enumerate}
\item \textbf{Matrix closure.}  Products of the displayed generators are
stored until no new matrix appears.  Projective closures normalize a single
nonzero entry only when the coefficient being tested is phase neutral.
\item \textbf{Direct determinants.}  Every finite represented matrix is
inserted into \eqref{nd:braid_tqft:eq:finite-def}.
\item \textbf{Power characters.}  Matrix powers and \eqref{nd:braid_tqft:eq:newton} compute
coefficients without diagonalizing; the numerical whole-function check uses
\eqref{nd:braid_tqft:eq:element-power} through power $60$.
\item \textbf{Monomial cycles.}  Nonzero column entries determine permutation
cycles and accumulated phases, independently checking
\eqref{nd:braid_tqft:eq:monomial-det}.
\item \textbf{Finite-field tracking.}  Weil words are followed simultaneously
in $SL_2(\F_3)$; affine fixed points are obtained from direct permutations and
from \eqref{nd:braid_tqft:eq:affine-fixed}.
\item \textbf{Orbit/design comparators.}  Pair orbits and the hook-length Haar
frame potential avoid the determinant formula.
\item \textbf{Haar diagnostics.}  Seeded continuous samples include standard
errors and group residuals and are explicitly reported as numerical evidence.
\end{enumerate}

Run the regression test from the repository root:
\begin{verbatim}
.venv/bin/python -m pytest -q \
  new_dists/braid_tqft_quantum_images_sigma_mgfs/test_verification.py
\end{verbatim}
Open the interactive laboratory with:
\begin{verbatim}
.venv/bin/marimo run \
  marimo_notebooks/braid_tqft.py
\end{verbatim}
The suite reports \textcolor{teal}{\textbf{PASS}}: 35 finite comparisons, 1,440 affine fixed-power
checks, and 16 compact-closure diagnostics.

\begin{thebibliography}{9}
\bibitem{nd:braid_tqft:howe}
S.~Howe, \emph{Random matrix statistics and zeroes of $L$-functions via
probability in $\lambda$-rings},
\url{https://arxiv.org/abs/2412.19295}.

\bibitem{nd:braid_tqft:flw}
M.~H.~Freedman, M.~J.~Larsen, and Z.~Wang,
\emph{The two-eigenvalue problem and density of Jones representation of braid
groups}, \url{https://arxiv.org/abs/math/0103200}.

\bibitem{nd:braid_tqft:thomas}
T.~Thomas, \emph{The Character of the Weil Representation},
\url{https://arxiv.org/abs/math/0610644}.

\bibitem{nd:braid_tqft:clifford4}
H.~Zhu, R.~Kueng, M.~Grassl, and D.~Gross,
\emph{The Clifford group fails gracefully to be a unitary 4-design},
\url{https://arxiv.org/abs/1609.08172}.

\bibitem{nd:braid_tqft:zhu-sic}
H.~Zhu, \emph{SIC--POVMs and Clifford groups in prime dimensions},
\url{https://arxiv.org/abs/1003.3591}.
\end{thebibliography}
